\section{Introducción: por qué este episodio es una clase sobre la trayectoria técnica multimodal}

Esta sección establece primero un marco de lectura. Este episodio es una entrevista de alta densidad técnica, poco común en el programa de entrevistas de negocios de Zhang Xiaojun (张小珺). Li Guangmi (李广密), como co-host, centra las preguntas en diez años de historia multimodal, las diferencias entre visión y lenguaje, lo que el paradigma de o1 sugiere para lo multimodal y el momento GPT-4 que podría llegar en los próximos dos años. El relato de Zhang Xiangyu (张祥雨) no es una colección de anécdotas sobre artículos, sino una trayectoria técnica: la CV aprendió de manera continua del NLP, pero la imagen estática en sí es distinta del lenguaje natural; el verdadero avance multimodal quizá requiera volver a entender next token prediction, RL, CoT, el espacio de acciones y el aprendizaje en línea.

\begin{importantbox}{Tesis central del episodio}
La dificultad de lo multimodal no se reduce a «poner imágenes, video y texto en un mismo modelo». En la imagen estática, la generación, la comprensión y la alineación con los humanos están separadas por naturaleza; la generación visual y la comprensión visual llevan mucho tiempo sin fusionarse realmente; después de o1, los investigadores empezaron a advertir que lo multimodal quizá también necesite cadenas de pensamiento, reflexión, ampliación del espacio de acciones y retroalimentación del entorno para que surja un verdadero momento GPT-4.
\end{importantbox}

\begin{knowledgebox}{Nota sobre la estrategia visual}
Este video es una entrevista con encuadre fijo, sin slides ni demostraciones en pantalla. El texto solo usa la portada para identificar la fuente; todas las imágenes del cuerpo son diagramas conceptuales de elaboración propia, que explican la historia de la transferencia de CV a NLP, la triple separación de la imagen estática, la gap entre generación y comprensión, los defectos de next token, RL/CoT, la Long CoT visual, la colaboración de varios modelos con long context y el aprendizaje en línea.
\end{knowledgebox}

\begin{figure}[H]
\centering
\includegraphics[width=0.96\textwidth]{cv-to-nlp-learning-history.png}
\caption{Historia de cómo la CV aprendió del NLP: de ImageNet/ResNet a ViT, iGPT, MIM y los grandes modelos multimodales. Diagrama conceptual de elaboración propia, basado en la conversación de 00:02:00--00:17:14.}
\end{figure}

\begin{knowledgebox}{Lectura de la figura: el momento GPT de la CV no copió simplemente al NLP}
La figura va de ImageNet y ResNet a ViT/iGPT/MIM y luego a los VLM, y muestra el proceso por el que la CV tomó una y otra vez paradigmas del NLP. La cuestión central es que transferir la arquitectura es fácil, pero transferir las propiedades de los datos y la función objetivo es muy difícil. Los datos de lenguaje provienen por naturaleza de la expresión humana, mientras que la imagen estática no contiene por naturaleza la comprensión humana.
\end{knowledgebox}

\subsection{Ruta de lectura}

Para no quedar abrumados por la terminología, esta sección plantea cuatro preguntas para leer el episodio. Primera: ¿por qué la CV aprendió del NLP durante diez años y aun así no reprodujo GPT con imágenes estáticas? Segunda: ¿por qué la generación y la comprensión de imágenes, aun dentro de un mismo sistema, siguen comportándose como dos modelos? Tercera: ¿por qué, cuanto más grande es el modelo, más fuertes son sus capacidades humanísticas y de conversación, pero su razonamiento matemático puede saltarse pasos y degradarse? Cuarta: ¿por qué la reflexión y los CoT pattern de o1 hicieron que la investigación multimodal volviera a ver un camino?

\noindent\begin{tabularx}{\textwidth}{p{0.19\textwidth}p{0.34\textwidth}X}
\toprule
Pregunta de lectura & Material de la entrevista & Juicio que hay que formarse \\
\midrule
Por qué a la CV le cuesta reproducir GPT & Imagen estática; separación entre generación, comprensión y alineación & La forma de los datos determina el objetivo de aprendizaje; no basta con transferir la arquitectura. \\
Por qué lo multimodal no se ha fusionado & Los modelos de comprensión y de generación mejoran cada uno por su lado, pero no se potencian entre sí & Falta un proceso de pensamiento común y un espacio de acciones controlable. \\
Por qué se degrada el razonamiento & Saltos de pasos en matemáticas de los modelos grandes; defectos de compresión de next token & Comprimir una distribución no equivale a precisión de cálculo. \\
Dónde está el próximo avance & RL, CoT, Long CoT visual, aprendizaje en línea & Lo decisivo está en el espacio de acciones, la retroalimentación y el aprendizaje continuo. \\
\bottomrule
\end{tabularx}

\subsection{Resumen de la sección}

El hilo conductor de EP102 es que lo multimodal no se reduce a «agregar un tipo más de entrada». Para entender el momento GPT-4 de los próximos dos años, hay que comprender a la vez la particularidad de los datos de imagen estática, los defectos de la función objetivo de los modelos de lenguaje, el modo de reflexión que introdujo o1 y el significado del aprendizaje en línea para la próxima generación de Agent.
